\section{Reproducibility and data}\label{app:repro}

\noindent The accompanying package contains the \LaTeX{} source, the bibliography, all scripts, machine-readable results,
and figures. \texttt{code/run\_all.sh} executes, in order, the experiments (\texttt{exp\_core.py}, \texttt{exp\_rates.py},
\texttt{exp\_power.py}, \texttt{exp\_contest.py}, \texttt{exp\_polymath.py}, \texttt{exp\_verify.py}, \texttt{exp\_speed.py},
\texttt{exp\_gss.py}, \texttt{exp\_village.py}, \texttt{exp\_cultures.py}, \texttt{exp\_world.py}, \texttt{exp\_notes.py}), the audit
(\texttt{audit.py}), the figures (\texttt{make\_figures.py}) and the manuscript inputs (\texttt{make\_tex.py}); total runtime
on one CPU core is about twenty minutes, most of it the evolutionary simulations and the bootstraps. Each script writes JSON
outputs to \texttt{results/} with its seed, sizes, parameters, and verification statistics. Figures are generated only from
those files and from the datasets. Every analysis-derived number in the text and every generated table is written
by \texttt{make\_tex.py} into \texttt{paper/generated/}, so the manuscript cannot drift from the results. The bibliography
is produced by \texttt{build\_bib.py} from Crossref records; the verification log is \texttt{results/ref\_verification.json}.
Of \RefCited{} references cited, \RefResolved{} have a DOI that was resolved against Crossref; abstracts could be retrieved
for most of them and specific claims were checked against those abstracts. Where no abstract was available the citation is
limited to what the title states. The manuscript compiles with \texttt{pdflatex}, \texttt{bibtex}, and two further
\texttt{pdflatex} passes.

\emph{Third-party data.} The fossil reanalysis uses three open datasets, redistributed in \texttt{data/} with attribution
and unmodified: Supplementary Dataset~1 of Gingerich~\cite{gingerich2022} (233 specimens; CC~BY~4.0); Supplementary File~2 of
Will et al.~\cite{will2021} (166 \Homo{} specimens with taxon, age and site coordinates; CC~BY~4.0); and the dataset of Du et
al.~\cite{du2018} (94 specimens; Dryad \texttt{doi:10.5061/dryad.c30g9}; CC0). Regional assignment of Middle Pleistocene
specimens uses site coordinates and is coded in \texttt{load\_fossils.py}. The present-day analysis uses the public
cumulative file of the General Social Survey, 1972--2024~\cite{gss2024}, as packaged for R by Healy~\cite{healygssr}. The file
is not redistributed: \texttt{code/fetch\_gss.py} downloads it and extracts the variables used, and \texttt{rda\_reader.py}
reads it without R. The speed-dating analysis uses the public file of the Columbia experiment~\cite{fisman2006}; it too is
not redistributed, and \texttt{code/fetch\_speed.py} copies it from a public mirror and checks its SHA-256 digest. Both
analyses use respondent-level public data with no identifiers, report only aggregates by sex and cohort, and make no
comparison among racial, ethnic or national groups. Three further sources are fetched by scripts and not redistributed:
the D-PLACE release of the ethnographic codes (\texttt{fetch\_dplace.py}, pinned to a commit), and the 175-country and
45-country survey files (\texttt{fetch\_open.py}, checked by SHA-256 digest; the first is licensed CC~BY~4.0). In those
analyses countries and societies are units of replication; none is named or ranked, and variables describing religion or
ethnicity are not read. No new data were collected from human
participants; all simulation data are synthetic.
